\section{World models and representations must face action tests}\label{sec:C34:05}
Predictive representations can improve perceptual and temporal structure while still omitting distinctions necessary for action. A visually coherent rollout can be wrong about a tool's side effects; a compact latent can predict typical observations while losing a rare precondition. The relevant research question is which representation preserves consequences under interventions, not merely which representation supports a strong passive probe.

Use controlled environments with known interventions to compare reconstructive, contrastive, JEPA-style, and action-conditioned representations at matched data and compute. Measure downstream planning or control with both frozen and adapted encoders. Include actions outside the behavior distribution and require uncertainty or abstention rather than extrapolation by default. Finally, validate selected plans through actual environment execution. A simulator-only gain can result from exploiting model error.

For tool environments, a world model can represent operation preconditions, possible effects, durations, and failure states rather than generate pixels. This typed model should be compared with direct execution, a learned unstructured predictor, and an oracle simulator in a controlled test bed. Replacing only the predictor with an oracle while holding the planner, representation interface, query budget, and evaluation population fixed estimates the effect of predictive error within that system; it is not a context-free decomposition of all errors. Comparing oracle-assisted planning with reactive control estimates a joint policy-and-information advantage under the declared budgets. Without these controls, a negative result cannot tell whether the model, planner, or task simply made planning unnecessary.

Generated imagery has no special evidential status. It can explain a mechanism or propose a possible future observation, but it cannot supply measured results. Diffusion and flow-based systems used as simulators need temporal, physical, and task-specific validation. Representation quality and generative realism are intermediate properties. Their practical importance depends on the decisions and failure events they change.

